\documentclass[lettersize,journal]{IEEEtran}
\usepackage{amsmath,amsfonts}
\usepackage{algorithmic}
\usepackage{algorithm}
\usepackage{array}
\usepackage[caption=false,font=normalsize,labelfont=sf,textfont=sf]{subfig}
\usepackage{textcomp}
\usepackage{stfloats}
\usepackage{url}
\usepackage{verbatim}
\usepackage{graphicx}
\usepackage{cite}
\hyphenation{op-tical net-works semi-conduc-tor IEEE-Xplore}
% updated with editorial comments 8/9/2021

% ═══════════════════════ PACKAGES ═══════════════════════
\usepackage[T1]{fontenc}
\usepackage[utf8]{inputenc}
\usepackage{amsmath,amssymb,amsfonts}
\usepackage{bm}
\usepackage{graphicx}
\usepackage{cite}
\usepackage{algorithmic}
\usepackage{algorithm}
\usepackage{url}
\usepackage{hyperref}
\usepackage{color}
\usepackage{booktabs}
\usepackage{multirow}
\usepackage{subcaption}
\usepackage{siunitx}
\usepackage{tikz}
\usepackage{stfloats}
\usetikzlibrary{arrows.meta,positioning,calc}

% ═══════════════════════ MATH COMMANDS ═══════════════════
\newcommand{\dq}{\hat{\bm{q}}}           % dual quaternion
\newcommand{\dqd}{\hat{\bm{q}}_d}        % desired dual quaternion
\newcommand{\dqe}{\hat{\bm{q}}_e}        % error dual quaternion
\newcommand{\se}{\mathfrak{se}(3)}        % Lie algebra se(3)
\newcommand{\so}{\mathfrak{so}(3)}        % Lie algebra so(3)
\newcommand{\SE}{SE(3)}                   % Special Euclidean group
\newcommand{\SO}{SO(3)}                   % Special Orthogonal group
\newcommand{\Spin}{\mathrm{Spin}(3)}      % Spin group
\newcommand{\R}{\mathbb{R}}              % Real numbers
\newcommand{\Hd}{\mathbb{H}_d}           % Dual quaternion space
\newcommand{\norm}[1]{\left\| #1 \right\|}
\DeclareMathOperator{\Log}{Log}
\DeclareMathOperator{\Exp}{Exp}
\DeclareMathOperator{\Ad}{Ad}
\DeclareMathOperator{\ad}{ad}
\DeclareMathOperator{\atan}{atan2}

%% Required preamble macros:

   \newcommand{\vect}[1]{\bm{#1}}
   \newcommand{\mat}[1]{\mathbf{#1}}

   \newcommand{\dt}{\Delta t}
   \newcommand{\st}{\text{s.t.}}
%% Required packages: amsmath, amssymb, bm, booktabs

% ═══════════════════════════════════════════════════════════

\begin{document}

%\title{Dual Quaternion Model Predictive Contouring Control with Lie-Invariant Error for Quadrotors}

\title{Lie-Invariant Dual Quaternion MPCC for Quadrotors}

\author{Bryan~S.~Guevara,
        Luis~F.~Simancas,
        and~Tiago~Nascimento

\thanks{Manuscript received January 14, 2025; revised March 3, 2025.}%
\thanks{Bryan~S.~Guevara is with the Laboratório de Automação
        e Sistemas de Robótica (LaSER), Universidade Federal da
        Paraíba (UFPB), João Pessoa, PB 58051-900, Brazil.
        {\tt\small bguevara@laser.ufpb.br}}%
\thanks{Luis~F.~Simancas is with Instituto de Automática (INAUT),
        Universidad Nacional de San Juan--CONICET, Argentina.
        {\tt\small lsimancas@inaut.unsj.edu.ar}}%
\thanks{Tiago~Nascimento is with the Laboratório de Automação
        e Sistemas de Robótica (LaSER), Universidade Federal da
        Paraíba (UFPB), João Pessoa, PB 58051-900, Brazil.
        {\tt\small tiago@laser.ufpb.br}}%
\thanks{This work has been supported by the Ministry of Science and Innovation of the Spanish Government through the research project PID2021-122685OBI00 and the grant PRE2019-088069 for the training of research PhD staff.}% <-this % stops a space

}

\markboth{IEEE Robotics and Automation Letters}%
{Guevara \MakeLowercase{\textit{et al.}}:
 Dual Quaternion MPCC with Lie-Invariant Error for Quadrotor}

\IEEEpubid{0000--0000/00\$00.00~\copyright~2025 IEEE}

\maketitle

\begin{abstract}
This paper presents a Model Predictive Contouring Control (MPCC)
framework for quadrotor UAVs in which the path-following error is
formulated directly on the Lie group $\SE$ using unit dual
quaternions. Unlike conventional MPCC approaches that decompose
position and orientation errors independently in Euclidean space,
the proposed method computes lag and contouring error components
through the logarithmic map of the dual quaternion error, yielding
a geometrically consistent, singularity-free decomposition in the
Lie algebra $\se$. The translational error is corrected via the
left Jacobian of $\SO$ to account for the rotation-translation
coupling intrinsic to $\se$, and projected onto the trajectory
tangent to separate lag and contouring terms under the
$\mathrm{Ad}$-invariant metric. The arc-length progress is
embedded as a system state and augmented with second-order virtual
progress dynamics, enabling the solver to jointly optimize the
physical control inputs and the progress acceleration within a
single nonlinear program. This preserves a torque-driven quadrotor
model while avoiding unrealistically abrupt changes in the virtual
traversal speed. The resulting OCP is
solved in real time via Sequential Quadratic Programming with
Real-Time Iteration (SQP-RTI) using the ACADOS framework. In
Model-in-the-Loop simulations over an \SI{80}{\meter} 3D Lissajous
trajectory, the proposed DQ-MPCC consistently reduces contouring and
orientation error relative to a Euclidean MPCC baseline, while exposing
a clear trade-off in lag error that is analyzed explicitly. These
results support the proposed formulation as a geometrically consistent
alternative when rigid-body pose accuracy is prioritized together with
real-time feasibility.
\end{abstract}

\begin{IEEEkeywords}
Model Predictive Contouring Control, Lie Groups, Dual Quaternions, Quadrotor UAV, $\SE$, Nonlinear MPC, Arc-Length Parametrization.
\end{IEEEkeywords}

%\IEEEpubidadjcol
\section{Introduction}
\IEEEPARstart{H}{igh-speed} autonomous flight requires controllers that
simultaneously maximize traversal speed and maintain precise spatial
accuracy along a prescribed path, a regime that is well known to demand
large attitude excursions and strongly coupled translational--rotational
motion in quadrotor systems~\cite{lee2010geometric,mellinger2011minimum}. Model Predictive Contouring Control
(MPCC)~\cite{liniger2015optimization} addresses this challenge by
decoupling \emph{spatial} tracking accuracy from \emph{temporal}
progress: rather than enforcing a fixed time schedule, the controller
treats the arc-length progress as an additional optimization variable,
enabling autonomous speed modulation—accelerating on straight segments
and decelerating in high-curvature regions. Since its introduction in
autonomous racing~\cite{liniger2015optimization}, MPCC has been
extended to quadrotor flight~\cite{romero2022model}, ground
vehicles~\cite{kabzan2019amz}, and multi-agent
systems, demonstrating consistent advantages over time-parametrized MPC
in agile, constraint-aware motion.

Despite these advances, virtually all existing MPCC formulations
decompose the path-following error in Euclidean space $\R^3$, relying
on Euler angles or rotation matrices to handle the orientational
degrees of freedom~\cite{romero2022model, kabzan2019amz}. These
representations carry well-documented drawbacks. Euler angles suffer
from gimbal lock and exhibit discontinuities at $\pm\pi$ that corrupt
gradient information inside the nonlinear program. Rotation matrices
require six orthonormality constraints, inflating the problem size and
complicating sensitivity analysis. More fundamentally, treating
position and orientation as independent quantities discards the
coupled rigid-body structure of $\SE$: any error metric defined
separately on $\R^3 \times SO(3)$ fails to capture the interaction
between rotational and translational motion that is intrinsic to
the Special Euclidean group.

Unit dual quaternions offer a principled remedy. A unit dual
quaternion $\dq \in \Hd$ encodes a full $\SE$ pose in a unified
8-dimensional algebra subject to only two scalar
constraints~\cite{wang2012, kenwright2012dual}, compared to twelve
for homogeneous matrices. Beyond compactness, dual quaternions inherit
the Lie group structure of $\SE$: the exponential and logarithmic maps
are closed-form, the error between two poses is a well-defined element
of the Lie algebra $\se$, and the associated $\mathrm{Ad}$-invariant
metric provides a geometrically consistent notion of distance on the
manifold. Control designs based on these properties have demonstrated
singularity-free tracking of coupled rotational-translational
motion~\cite{wang2012,abaunza2017,johansen2019review}, yet no prior work has
exploited them within an MPCC framework for path-following speed
optimization.

\IEEEpubidadjcol

This paper closes that gap. We present a Lie-invariant MPCC framework
for quadrotor UAVs in which the lag and contouring error components
are computed directly in $\se$ via the logarithmic map of the dual
quaternion pose error. The arc-length progress is embedded as a system
state and augmented with second-order virtual-progress dynamics
\(
\dot{s} = v_s,\;
\dot{v}_s = a_s
\),
allowing the solver to jointly optimize physical control inputs and the
progress acceleration within a single nonlinear program. This choice is
inspired by the progress regularisation strategy in~\cite{romero2022model},
but adapted here to a torque-driven MPCC formulation so that the baseline
actuation structure is preserved. The resulting optimal control problem is
solved in real time using Sequential Quadratic Programming with
Real-Time Iteration (SQP-RTI) inside the ACADOS
framework~\cite{verschueren2021acados}.

The main contributions of this work are as follows:
\begin{enumerate}
    \item \textbf{Dual-quaternion MPCC on $\SE$:} To the best of our
    knowledge, we present the first MPCC formulation for quadrotors
    that operates directly on the Lie group $\SE$ using unit dual
    quaternions, extending contouring control from Euclidean path
    following to full rigid-body motion with coupled rotation and
    translation.

    \item \textbf{Lie-algebra lag-contouring decomposition:} We
    introduce a geometrically consistent decomposition of the
    path-following error in the Lie algebra $\se$. The pose error is
    obtained from the dual quaternion logarithmic map; its
    translational component is corrected via the left Jacobian of
    $\SO$ and projected onto the trajectory tangent to yield
    coordinate-invariant lag and contouring terms consistent with the
    $\mathrm{Ad}$-invariant metric of $\se$.

    \item \textbf{Second-order virtual-progress augmentation:} We
    augment the MPCC progress variable with explicit second-order
    dynamics, introducing a bounded progress acceleration as the
    virtual optimization input. In contrast to actuator-level
    formulations such as~\cite{romero2022model}, this preserves the
    thrust--torque structure of the quadrotor model used in both the
    Euclidean baseline and the proposed DQ-MPCC, while still enforcing
    smooth and dynamically plausible progress evolution.

    \item \textbf{Real-time validation:} We demonstrate real-time
    feasibility at \SI{100}{\hertz} on an \SI{80}{\meter} 3D Lissajous
    path using SQP-RTI, and show through MiL experiments that the
    proposed formulation improves contouring and orientation accuracy
    while preserving comparable traversal times with respect to the
    Euclidean baseline.
\end{enumerate}

The remainder of the paper is organized as follows.
Section~\ref{sec:prelim} reviews dual quaternion algebra and the
relevant Lie group structure. Section~\ref{sec:dynamics} derives the
UAV dynamics in dual quaternion form. Section~\ref{sec:mpcc} presents the Lie-invariant MPCC
formulation. Section~\ref{sec:impl} details the implementation and
solver configuration. Section~\ref{sec:results} reports simulation
results, and Section~\ref{sec:conclusion} concludes with directions
for future work.
% ═══════════════════════ II. PRELIMINARIES ══════════════════
\section{Mathematical Preliminaries}
\label{sec:prelim}



\subsection{Dual Quaternion Algebra}

A dual number is an expression $a + \varepsilon b$, where $a, b \in \mathbb{R}$
and the dual unit $\varepsilon$ satisfies $\varepsilon^{2} = 0$,
$\varepsilon \neq 0$. Extending this to the quaternion algebra
$\mathbb{H}$, a \emph{dual quaternion} is
\begin{equation}
    \dq = q_r + \varepsilon\, q_d, \qquad q_r,\, q_d \in \mathbb{H},
\end{equation}
where $q_r$ is the \emph{real part} and $q_d$ is the \emph{dual part}.
This representation is standard in rigid-body kinematics and provides a
compact pose parametrization that remains closed under composition and
inversion~\cite{daniilidis1999,selig2005}.
The product of two dual quaternions follows from the Hamilton product
$\otimes$ applied to both parts:
\begin{equation}
    \dq_1 \otimes \dq_2
        = q_{r_1} \otimes q_{r_2}
        + \varepsilon\!\left(
            q_{r_1} \otimes q_{d_2} + q_{d_1} \otimes q_{r_2}
          \right).
    \label{eq:dq_product}
\end{equation}

A dual quaternion represents a rigid-body pose in $\SE$ when it
satisfies the \emph{unit constraints}
\begin{equation}
    \|q_r\| = 1, \qquad \langle q_r,\, q_d \rangle = 0,
    \label{eq:unit_constraints}
\end{equation}
where $\langle\cdot,\cdot\rangle$ denotes the scalar (dot) product of
two quaternions treated as vectors in $\mathbb{R}^{4}$.
These conditions characterize the unit dual quaternions associated with
rigid motions and ensure consistency with the geometry of $\SE$~\cite{selig2005}.
Given a rotation parametrized by angle $\theta$ about unit axis
$\hat{\bm{n}} \in \mathbb{R}^{3}$ and a translation
$\bm{t} \in \mathbb{R}^{3}$, the corresponding unit dual quaternion is
\begin{equation}
    \dq = q_r + \frac{\varepsilon}{2}\, q_t \otimes q_r,
    \qquad
    q_r = \begin{bmatrix}\cos\tfrac{\theta}{2}\\[2pt]
                          \hat{\bm{n}}\sin\tfrac{\theta}{2}\end{bmatrix},
    \quad
    q_t = \begin{bmatrix}0\\ \bm{t}\end{bmatrix}.
    \label{eq:dq_pose}
\end{equation}
The conjugate of a unit dual quaternion is
$\dq^{*} = q_r^{*} + \varepsilon\, q_d^{*}$, and the inverse equals
the conjugate, $\dq^{-1} = \dq^{*}$, by virtue of the unit
constraint~\eqref{eq:unit_constraints}. The translation and rotation
are recovered from any unit dual quaternion as
\begin{equation}
    \bm{t} = 2\,\bigl(q_d \otimes q_r^{*}\bigr)\big|_{\mathrm{vec}},
    \qquad
    R = R(q_r) \in SO(3),
    \label{eq:dq_extract}
\end{equation}
where $(\cdot)|_{\mathrm{vec}}$ denotes the vector part of a quaternion.



\begin{figure}[t]
    \centering
    \includegraphics[width=\columnwidth]{DQ-MPCC.png}
    \caption{
Geometric interpretation of the Lie-invariant lag--contouring
decomposition on $\SE$.
The desired trajectory $\bm{\gamma}(s)$ defines a curve on $\SE$,
parameterized by the arc length $s$.
At time step $k$, the reference and current poses are denoted by
$q_k^d$ and $q_k$, respectively, with relative error
$q^e = q_k^{d*}\otimes q_k \in \SE$.
Applying the logarithmic map yields the Lie-algebra error
$\bm{\xi}^e=(\bm{\phi}^e,\bm{\rho})\in\mathfrak{se}(3)$ in the tangent
space at $q_k^d$.
Here, $\bm{\phi}^e\in\mathfrak{so}(3)$ is the rotational error, whereas
$\bm{\rho}\in\mathbb{R}^3$ is the translational component of the
logarithmic pose error, and thus is not the Euclidean position
difference, since it is coupled to the relative orientation.
Using the unit path tangent $\vect{e}_t^d(s)=d\bm{p}^d/ds$,
$\bm{\rho}$ is decomposed into a lag component
$\bm{\rho}_{\mathrm{lag}}$, aligned with the path, and a contouring
component $\bm{\rho}_{\mathrm{cont}}$, orthogonal to it.
This yields a geometrically consistent separation of longitudinal and
lateral path-following errors in $\mathfrak{se}(3)$.
    }
    \label{fig:lie_decomposition}
\end{figure}

\subsection{Lie Group Structure and Logarithmic Map}

The set of unit dual quaternions satisfying~\eqref{eq:unit_constraints}
forms the group $\mathrm{Spin}(3) \ltimes \mathbb{R}^{3}$, the double
cover of $\SE$: every rigid-body pose corresponds to exactly two
antipodal unit dual quaternions $\pm\dq$~\cite{kenwright2012dual,selig2005}.
Its Lie algebra $\se$ consists of six-dimensional vectors
\begin{equation}
    \bm{\xi} = \begin{pmatrix}\bm{\phi}\\ \bm{\rho}\end{pmatrix} \in \mathbb{R}^{6},
\end{equation}
where $\bm{\phi} \in \mathbb{R}^{3}$ is the rotation vector and
$\bm{\rho} \in \mathbb{R}^{3}$ is the coupled translational component.

The \textbf{logarithmic map} $\Log: \Hd \to \se$ proceeds in two steps.
First, the rotation angle and axis are extracted from the real part $q_r$:
\begin{equation}
    \theta = 2\operatorname{atan2}\!\bigl(\|q_{r,\mathrm{vec}}\|,\;q_{r,0}\bigr),
    \qquad
    \hat{\bm{n}} = \frac{q_{r,\mathrm{vec}}}{\|q_{r,\mathrm{vec}}\|},
    \label{eq:log_rotation}
\end{equation}
and the translation $\bm{t}$ is recovered via~\eqref{eq:dq_extract}.
Then the Lie algebra element is
\begin{equation}
    \Log(\dq)
    = \begin{pmatrix}\bm{\phi}\\ \bm{\rho}\end{pmatrix}
    = \begin{pmatrix}\theta\,\hat{\bm{n}} \\[3pt] \bm{t}\end{pmatrix}.
    \label{eq:log_map}
\end{equation}
Throughout the paper, quantities associated with the desired
(reference) pose carry a superscript $d$: $\dq^{d}$, $q_r^{d}$,
$\bm{p}^{d}$, $R^{d}$, etc.
Hemisphere consistency is enforced by
\begin{equation}
    q_r \;\leftarrow\; \mathrm{sign}(q_{r,0})\,q_r
    \label{eq:hemisphere}
\end{equation}
before evaluating~\eqref{eq:log_rotation}--\eqref{eq:log_map}.
These logarithmic coordinates follow the standard Lie-theoretic local
representation used in robotics, where Jacobian corrections map finite
group displacements to tangent perturbations~\cite{sola2018micro,barfoot2017}.




\section{Dual Quaternion Dynamics of the Quadrotor}
\label{sec:dynamics}

\subsection{Classical Quadrotor Model}

Consider a quadrotor of mass $m$ and inertia matrix
$\mat{J} \in \mathbb{R}^{3\times 3}$. Let
$\vect{p} \in \mathbb{R}^{3}$ denote the position in the inertial
frame, $q_r \in \mathbb{S}^{3}$ the unit attitude quaternion,
$\vect{v} \in \mathbb{R}^{3}$ the translational velocity in the
inertial frame, and $\vect{\omega} \in \mathbb{R}^{3}$ the angular
velocity in the body frame. The equations of motion are:
\begin{align}
    \dot{\vect{p}}      &= \vect{v}, \label{eq:uav_p}\\
    \dot{q}_r           &= \tfrac{1}{2}\, q_r \otimes \vect{\omega},
                                                  \label{eq:uav_q}\\
    \dot{\vect{v}}      &= \frac{f}{m}\, R(q_r)\,\vect{e}_3 + \vect{g},
                                                  \label{eq:uav_v}\\
    \dot{\vect{\omega}} &= \mat{J}^{-1}\!\bigl(
                            \vect{\tau} - \vect{\omega}
                            \times (\mat{J}\vect{\omega})
                          \bigr),                \label{eq:uav_w}
\end{align}
where $f \in \mathbb{R}$ is the total thrust,
$\vect{\tau} \in \mathbb{R}^{3}$ is the control torque vector,
$\vect{e}_3 = [0,0,1]^\top$ is the body-frame thrust direction,
and $\vect{g} = [0,0,-g]^\top$\,m/s$^2$ is the gravitational
acceleration in the inertial frame.

\subsection{Dual Quaternion State Representation}

The separate kinematics~\eqref{eq:uav_p}--\eqref{eq:uav_q} are
unified into a single pose equation by assembling the unit dual
quaternion~\eqref{eq:dq_pose}:
\begin{equation}
    \dq = q_r + \varepsilon\, q_d,
    \label{eq:dq_parts}
\end{equation}
where $q_r \in \mathbb{R}^{4}$ is the real part encoding the
attitude and $q_d \in \mathbb{R}^{4}$ is the dual part encoding the
translation coupled with the rotation, defined as
\begin{equation}
    q_d \;\triangleq\; \tfrac{1}{2}\,q_t \otimes q_r,
    \qquad \text{with} \qquad q_t = [0,\;\vect{p}^\top]^\top.
    \label{eq:qd_def}
\end{equation}

The angular and linear velocities are packed into the
\emph{dual twist}, a pure dual quaternion
\begin{equation}
    \hat{\vect{\omega}} = \vect{\omega} + \varepsilon\,\vect{v}_b,
    \label{eq:dual_twist}
\end{equation}
where
\begin{equation}
    \vect{v}_b \;\triangleq\;
    q_r^{*} \otimes \vect{v} \otimes q_r
    \;=\; \mathrm{Ad}_{q_r^{*}}(\vect{v})
    \label{eq:vb_def}
\end{equation}
is the translational velocity rotated from the inertial frame into
the body frame via the adjoint action.
Because $\hat{\vect{\omega}}$ is pure (both scalar parts vanish),
it carries exactly six independent degrees of freedom collected in
$\vect{\xi} = [\vect{\omega}^\top,\;\vect{v}_b^\top]^\top
\in \mathbb{R}^{6}$.

With these definitions the pose kinematics reduce to
\begin{equation}
    \dot{\dq} = \tfrac{1}{2}\,\dq \otimes \hat{\vect{\omega}},
    \label{eq:dq_kinematics}
\end{equation}
and the dual twist dynamics decompose as
\begin{equation}
    \dot{\hat{\vect{\omega}}} = \hat{\vect{f}} + \hat{\vect{u}},
    \label{eq:dual_twist_dyn}
\end{equation}
where the drift and control dual quaternions are, respectively,
\begin{align}
    \hat{\vect{f}} &=
        -\mat{J}^{-1}\!\bigl(\vect{\omega}\times(\mat{J}\vect{\omega})\bigr)
        + \varepsilon\!\left(
            \vect{v}_b\times\vect{\omega}
            + \mathrm{Ad}_{q_r^{*}}(\vect{g})
          \right),  \label{eq:f_drift}\\[4pt]
    \hat{\vect{u}} &=
        \mat{J}^{-1}\vect{\tau}
        + \varepsilon\!\left(\frac{f}{m}\,\vect{e}_3\right).
    \label{eq:u_control}
\end{align}



\subsection{Compact State Equation}

Collecting the pose $\dq = q_r + \varepsilon\, q_d$ and the dual
twist $\hat{\vect{\omega}} = \vect{\omega} + \varepsilon\,\vect{v}_b$
into a single state, the complete UAV dynamics in dual quaternion
form are:
\begin{equation}
\boxed{
\frac{d}{dt}
\begin{bmatrix} \dq \\ \hat{\vect{\omega}} \end{bmatrix}
=
\begin{bmatrix}
    \tfrac{1}{2}\,\dq \otimes \hat{\vect{\omega}} \\[4pt]
    \hat{\vect{f}} + \hat{\vect{u}}
\end{bmatrix}.
}
\label{eq:dynamics}
\end{equation}
This compact form follows the usual Lie-group treatment of rigid-body
kinematics and dynamics, where pose evolution is written directly on the
motion group instead of splitting translation and attitude into unrelated
coordinate charts~\cite{park1995,murray1994mathematical}.

For numerical implementation, $\dq = q_r + \varepsilon\,q_d$ is
stored as its eight scalar components and
$\hat{\vect{\omega}} = \vect{\omega} + \varepsilon\,\vect{v}_b$
as its six vector components $\vect{\xi}$, yielding the state vector
\begin{equation}
    \vect{x} =
    \bigl[\,
        \underbrace{q_r \in \mathbb{R}^4}_{\text{attitude}},\;
        \underbrace{q_{d} \in \mathbb{R}^4}_{\text{dual part}},\;
        \underbrace{\vect{\omega} \in \mathbb{R}^3}_{\text{ang.\ vel.}},\;
        \underbrace{\vect{v}_b \in \mathbb{R}^3}_{\text{body lin.\ vel.}}
    \,\bigr]^\top \in \mathbb{R}^{14},
    \label{eq:state_vector}
\end{equation}
and the control input
$\vect{u} = [f,\;\tau_1,\;\tau_2,\;\tau_3]^\top \in \mathbb{R}^{4}$.
The unit-norm constraint $\|q_r\| = 1$ is maintained via a
Baumgarte-type correction $K(1-\|q_r\|)\,q_r$ with $K=10$ added
to $\dot{q}_r$ inside the integrator, supplemented by a soft
penalty constraint in the OCP formulation. This is a classical
constraint-stabilization device used here only to keep the numerical
integration close to the unit-quaternion manifold~\cite{baumgarte1972}.


\subsection{Reference Trajectory and Arc-Length Parametrization}
\label{sec:ref_traj}

MPCC requires a \emph{spatial} reference indexed by distance along
the path rather than reference path by time $\vect{p}^d(t)$, so that the optimizer can freely
choose how fast to traverse each segment. The arc length is
\begin{equation}
    s(t) = \int_0^t \bigl\|\dot{\vect{p}}^d(\tau)\bigr\|\,d\tau,
    \label{eq:arc_length}
\end{equation}
and the inverse $t(s)$ is recovered by cubic spline interpolation.
The reference pose at station $s$ is
\begin{equation}
    \bm{\gamma}(s)
        = \bigl(\vect{p}^d(s),\; q_r^d(s)\bigr) \in \SE,
    \label{eq:gamma}
\end{equation}
where the desired attitude $q_r^d(s)$ is constructed from the path
geometry. First, the unit tangent
\begin{equation}
    \vect{t}^d(s)
    = \frac{d\vect{p}^d/ds}{\|d\vect{p}^d/ds\|}
    \label{eq:heading_tangent}
\end{equation}
defines the desired heading direction. Next, the path curvature vector
is approximated by
\begin{equation}
    \vect{\kappa}^d(s) = \frac{d\vect{t}^d}{ds},
    \label{eq:curvature_vector}
\end{equation}
and combined with a nominal reference speed $\bar{v}_s$ to form the
feedforward lateral acceleration
\begin{equation}
    \vect{a}_{\mathrm{lat}}^d(s) = \bar{v}_s^2\,\vect{\kappa}^d(s).
    \label{eq:lateral_acc_ref}
\end{equation}
The desired thrust direction is then chosen as
\begin{equation}
    \vect{b}_3^d(s)
    = \frac{\vect{a}_{\mathrm{lat}}^d(s) + g\,\vect{e}_3}
           {\|\vect{a}_{\mathrm{lat}}^d(s) + g\,\vect{e}_3\|},
    \label{eq:b3_des}
\end{equation}
which is clipped to a maximum admissible tilt angle in order to keep
the reference physically realizable. The desired body x-axis is the
projection of the tangent onto the plane orthogonal to $\vect{b}_3^d$,
\begin{equation}
    \vect{b}_1^d(s)
    = \frac{
        \left(\bm{I}_3 - \vect{b}_3^d{\vect{b}_3^d}^{\!\top}\right)\vect{t}^d(s)
      }{
        \left\|
        \left(\bm{I}_3 - \vect{b}_3^d{\vect{b}_3^d}^{\!\top}\right)\vect{t}^d(s)
        \right\|
      },
    \label{eq:b1_des}
\end{equation}
with $\vect{b}_2^d(s)=\vect{b}_3^d(s)\times\vect{b}_1^d(s)$. The
reference rotation matrix is finally assembled as
\begin{equation}
    R^d(s) = \begin{bmatrix}\vect{b}_1^d(s) & \vect{b}_2^d(s) & \vect{b}_3^d(s)\end{bmatrix}
    \in SO(3),
    \label{eq:Rd}
\end{equation}
and mapped to the quaternion $q_r^d(s)$. Unlike a yaw-only reference,
this construction makes the orientation reference compatible with the
centripetal acceleration required for fast path following.


% ═══════════════════════ III. PROBLEM FORMULATION ══════════

\section{Lie-Invariant MPCC Formulation}
\label{sec:mpcc}

Classical MPCC decomposes the path-following error in Euclidean space
using the Frenet--Serret frame~\cite{liniger2015optimization}, thereby
treating position and orientation separately and disregarding the
underlying manifold structure of $\SE$. In contrast, the proposed
formulation performs the decomposition directly in $\se$ through the
logarithmic map of the dual-quaternion error, together with the
left-Jacobian correction introduced in
Section~\ref{sec:left_jac}.

\subsection{Baseline MPCC}
\label{sec:baseline}

To assess the benefit of the proposed geometric formulation, the
DQ-MPCC is benchmarked against a standard Euclidean MPCC baseline.
Given the reference position $\vect{p}^d(s)$ and its unit tangent
$\vect{t}(s)$, the position tracking error
$\vect{e}_t = \vect{p}^d(s) - \vect{p}$ is decomposed into lag and
contouring components as
\begin{equation}
    \vect{e}_\ell = \bigl(\vect{t}^\top \vect{e}_t\bigr)\,\vect{t},
    \qquad
    \vect{e}_c = \bigl(\bm{I}_3 - \vect{t}\,\vect{t}^\top\bigr)\,\vect{e}_t,
\end{equation}
while the attitude error is defined through the quaternion logarithm
\begin{equation}
    \bm{\eta} = \Log(q^{-1} \otimes q^d) \in \so,
\end{equation}
using the same curvature-compatible attitude reference
$q^d(s)=q_r^d(s)$ introduced in~\eqref{eq:b3_des}--\eqref{eq:Rd}. This
avoids penalizing a hover-like orientation while the position loop
demands the tilt required for aggressive path following.

Both controllers share the same solver configuration, prediction
horizon, stage-cost structure~\eqref{eq:cost}, and physical
constraints~\eqref{eq:ocp_thrust}--\eqref{eq:ocp_speed}. Hence, the
only difference lies in the pose-error representation: the baseline
uses a Euclidean decomposition in $\mathbb{R}^3$ with a separate
quaternion attitude error, whereas the proposed DQ-MPCC employs the
geometrically consistent logarithmic map on $\SE$, as detailed in
Section~\ref{sec:mpcc}. In the baseline cost, $\bm{\eta}$ replaces
$\bm{\phi}^e$ with the same quadratic penalty structure as
in~\eqref{eq:cost}. Therefore, any observed performance difference can
be attributed solely to the adopted error representation.







\subsection{DQ-MPCC Formulation}

Let $\dq^d(s) = \bm{\gamma}(s)$ denote the reference dual quaternion,
evaluated at the current arc-length station $s$ --- the fifteenth
component of the augmented state~\eqref{eq:aug_state}; the sixteenth
component stores the virtual progress speed $v_s$. The pose error is
defined as the left-invariant error on $\SE$:
\begin{equation}
    \dq^e = \dq^{d\,*}(s) \otimes \dq,
    \label{eq:dq_error}
\end{equation}
which expresses the deviation of the current pose $\dq$ from the
reference $\dq^d(s)$ in the body frame of the reference~\cite{bullo1995proportional}. After
applying hemisphere correction~\eqref{eq:hemisphere}, the
logarithmic map~\eqref{eq:log_map} yields the Lie algebra error
\begin{equation}
    \bm{\xi}^e = \Log\!\left(\dq^e\right)
               = \begin{pmatrix} \bm{\phi}^e \\ \bm{\rho}^e
                 \end{pmatrix} \in \se,
    \label{eq:log_error}
\end{equation}
where $\bm{\phi}^e \in \mathbb{R}^3$ is the rotational error vector
and $\bm{\rho}^e \in \mathbb{R}^3$ is the raw translational
component. As detailed in Section~\ref{sec:left_jac}, $\bm{\rho}^e$
encodes translation in a frame twisted by $\bm{\phi}^e$ through the
screw structure of $\se$; the corrected translational error
$\bm{\rho}$ is defined in~\eqref{eq:rho_corrected}.

%\subsection{Left Jacobian of $\SO$ and Translational Correction}
%\label{sec:left_jac}

A key quantity linking Lie-algebra errors to Euclidean position errors
is the \emph{left Jacobian} of $\SO$,
$\bm{J}_{l}(\bm{\phi}) \in \mathbb{R}^{3\times 3}$~\cite{sola2018micro}:
\begin{equation}
    \bm{J}_{l}(\bm{\phi})
    = \bm{I}_{3}
      + \frac{1 - \cos\theta}{\theta^{2}}\,[\bm{\phi}]_{\times}
      + \frac{\theta - \sin\theta}{\theta^{3}}\,
        [\bm{\phi}]_{\times}^{2},
    \label{eq:left_jacobian}
\end{equation}
where $\theta = \|\bm{\phi}\|$ and
$[\bm{\phi}]_{\times}$ is the $3\times 3$ skew-symmetric matrix
associated with $\bm{\phi}$.
Its inverse admits the closed-form expression~\cite{sola2018micro}:
\begin{equation}
    \bm{J}_{l}(\bm{\phi})^{-1}
    = \frac{\theta}{2}\cot\!\frac{\theta}{2}\;\bm{I}_{3}
      \;-\;\frac{\theta}{2}\,[\hat{\bm{n}}]_{\times}
      \;+\;\Bigl(1 - \frac{\theta}{2}\cot\!\frac{\theta}{2}\Bigr)
            \,\hat{\bm{n}}\,\hat{\bm{n}}^{\top},
    \label{eq:left_jacobian_inv}
\end{equation}
where $\hat{\bm{n}} = \bm{\phi}/\theta$ is the unit rotation axis.
Using the analytical form~\eqref{eq:left_jacobian_inv} instead of a
numerical matrix inversion avoids the heavier expression graph that
\texttt{ca.inv} would introduce in the CasADi symbolic framework,
yielding cleaner automatic differentiation and faster code generation
within the acados solver.
For $\theta < 10^{-7}$ both expressions reduce to
$\bm{J}_{l} \to \bm{I}_{3}$ and
$\bm{J}_{l}^{-1} \to \bm{I}_{3}$
to avoid numerical singularity.

To understand why this Jacobian is needed, recall that the
translational component $\bm{\rho}^e$ returned by $\Log(\dq^e)$
via~\eqref{eq:log_map} is \emph{not} the Euclidean displacement
$\bm{p} - \bm{p}^{d}$. Rather, it is the translation expressed in
a frame that has been \emph{twisted} by the rotation $\bm{\phi}^e$
through the $\se$ screw structure~\cite{sola2018micro}. The
corrected translational error is therefore defined as
\begin{equation}
    \bm{\rho} = \bm{J}_{l}(\bm{\phi}^e)^{-1}\, R^{d\,\top}
                \bigl(\bm{p} - \bm{p}^{d}\bigr),
    \label{eq:rho_corrected}
\end{equation}
where $R^{d} = R(q_r^{d}) \in SO(3)$ is the desired rotation matrix
and $\bm{p} - \bm{p}^{d} \in \mathbb{R}^{3}$ is the position error
in the inertial frame. Equation~\eqref{eq:rho_corrected} thus defines
a vector in $\mathbb{R}^{3}$ that (i) is expressed in the desired
body frame via $R^{d\,\top}$, and (ii) accounts for the
rotation-translation coupling of $\se$ via $\bm{J}_{l}^{-1}$.


%\subsection{Lag-Contouring Decomposition in $\se$}

The corrected translational error $\bm{\rho}$~\eqref{eq:rho_corrected}
is decomposed into components parallel and perpendicular to the
instantaneous tangent direction of the reference trajectory. At
arc-length station $s$, this tangent is
\begin{equation}
    \vect{e}^d_t(s)
        = \frac{\vect{p}^{d\prime}(s)}{\|\vect{p}^{d\prime}(s)\|}, \qquad
        \text{with} 
    \qquad
    \vect{p}^{d\prime}(s) = \frac{d\vect{p}^d}{ds},
    \label{eq:tangent}
\end{equation}
where $\vect{p}^{d\prime}(s)$ is the derivative of the reference
position with respect to arc length, available analytically from
the CasADi interpolant of Section~\ref{sec:ref_traj}. Because the
parametrization is by arc length, $\|\vect{p}^{d\prime}\| = 1$
holds exactly along smooth segments, so~\eqref{eq:tangent} is
well-defined everywhere except at cusps. The lag and contouring
components are then
\begin{subequations}
\label{eq:lag_contour}
\begin{align}
    \bm{\rho}_\mathrm{lag}
        &= \bigl(\vect{e}^d_t\,{\vect{e}^d_t}^{\!\top}\bigr)
           \bm{\rho},
    \label{eq:rho_lag}\\[2pt]
    \bm{\rho}_\mathrm{cont}
        &= \bigl(\bm{I}_3
           - \vect{e}^d_t\,{\vect{e}^d_t}^{\!\top}\bigr)
           \bm{\rho}.
    \label{eq:rho_cont}
\end{align}
\end{subequations}
The \emph{lag error} $\bm{\rho}_\mathrm{lag}$ measures how far the
vehicle lags behind or leads the virtual reference point along the
path; the \emph{contouring error} $\bm{\rho}_\mathrm{cont}$ measures
lateral deviation from it.

This decomposition is geometrically consistent with the $\SE$
structure. Because $\bm{\rho}$ is obtained via the left Jacobian
inverse~\eqref{eq:rho_corrected}, it resides in the translational
subspace of $\se$ with rotation-translation coupling accounted for.
Both $\bm{\rho}_\mathrm{lag}$ and $\bm{\rho}_\mathrm{cont}$ are
orthogonal projections under the left-invariant metric
\begin{equation}
    \langle \bm{\xi}_1, \bm{\xi}_2 \rangle_{\se}
        = \lambda\,\bm{\phi}_1^\top\bm{\phi}_2
        + \mu\,\bm{\rho}_1^\top\bm{\rho}_2,
    \qquad \lambda,\mu > 0,
    \label{eq:se3_metric}
\end{equation}
which defines the weighted inner product adopted in this work on the
tangent coordinates of $\se$~\cite{murray1994mathematical}. A purely Euclidean
decomposition of $\vect{p} - \vect{p}^d$ would not respect this
structure and would produce biased error signals whenever
$\bm{\phi}^e \neq \vect{0}$.

\subsection{Optimal Control Problem}

To embed arc-length progress as a degree of freedom while preventing
unrealistically abrupt changes in the virtual traversal speed, the
dynamics~\eqref{eq:dynamics} are augmented with the progress state
$s \in [0, s_{\max}]$, the virtual progress speed $v_s$, and the
virtual progress acceleration input $a_s$:
%
\begin{equation}
    \dot{s} = v_s,
    \qquad
    \dot{v}_s = a_s,
    \label{eq:arc_dyn}
\end{equation}
%
yielding the extended state vector
%
\begin{equation}
    \vect{x}_{\mathrm{a}}
    = \bigl[\,
        \underbrace{\vect{x} \in \mathbb{R}^{14}}_{\text{UAV state}},\;
        \underbrace{s \in \mathbb{R}}_{\text{arc length}},\;
        \underbrace{v_s \in \mathbb{R}}_{\text{virtual speed}}
    \,\bigr]^\top \in \mathbb{R}^{16},
    \label{eq:aug_state}
\end{equation}
%
and the augmented control vector
%
\begin{equation}
    \vect{u}_{\mathrm{a}}
    = \bigl[\,
        \underbrace{\vect{u}^\top \in \mathbb{R}^{4}}_{\text{physical inputs}},\;
        \underbrace{a_s \in \mathbb{R}}_{\text{virtual acceleration}}
    \,\bigr]^\top \in \mathbb{R}^{5},
    \label{eq:aug_control}
\end{equation}
%

The stage cost at each prediction node is:
\begin{equation}
\begin{aligned}
    \ell(\vect{x}_{\mathrm{a}}, \vect{u}_{\mathrm{a}})
        &= {\bm{\xi}^e}^\top \bm{Q}_\phi\,\bm{\xi}^e
         + \bm{\rho}_\mathrm{lag}^\top \bm{Q}_\ell\,\bm{\rho}_\mathrm{lag}
         + \bm{\rho}_\mathrm{cont}^\top \bm{Q}_c\,\bm{\rho}_\mathrm{cont}
           \\[4pt]
        &\quad
         + \vect{u}^\top \bm{W}_u\,\vect{u}
         + \vect{\omega}^\top \bm{Q}_\omega\,\vect{\omega}
         + Q_a\,a_s^2
         - Q_s\,v_s .
\end{aligned}
\label{eq:cost}
\end{equation}

To avoid double-penalising translation, $\bm{Q}_\phi \in \mathbb{R}^{6\times 6}$
is structured with upper block $\bm{Q}_q$ acting on the rotational
component $\bm{\phi}^e$ of $\bm{\xi}^e$ and a zero lower block,
since lateral and longitudinal errors are handled exclusively by
$\bm{Q}_c$ and $\bm{Q}_\ell$. All tracking and regularisation weights
$\bm{Q}_\ell,\,\bm{Q}_c,\,\bm{Q}_q,\,\bm{Q}_\omega \in \mathbb{R}^{3\times 3}$
are diagonal, allowing independent penalisation per axis and
accommodating the typical asymmetry between vertical and horizontal
dynamics of a quadrotor. The control effort weight
$\bm{W}_u
\in \mathbb{R}^{4\times 4}$ penalises thrust and torques independently,
and $Q_s > 0$ is a scalar speed incentive weight. Numerical values of
all weights are reported in Table~\ref{tab:cost_weights}. The scalar
$Q_a > 0$ regularises the virtual progress acceleration and is kept
fixed in the implementation to isolate the effect of the main MPCC
tracking weights during bilevel tuning.

The terminal cost retains all error terms but omits control
regularization:
\begin{equation}
    \ell_e(\vect{x}_{\mathrm{a}})
        = {\bm{\xi}^e}^\top \bm{Q}_\phi\,\bm{\xi}^e
        + \bm{\rho}_\mathrm{lag}^\top \bm{Q}_\ell\,\bm{\rho}_\mathrm{lag}
        + \bm{\rho}_\mathrm{cont}^\top \bm{Q}_c\,\bm{\rho}_\mathrm{cont}
        + \vect{\omega}^\top \bm{Q}_\omega\,\vect{\omega}.
    \label{eq:terminal_cost}
\end{equation}


The complete OCP of DQ-MPCC over a prediction horizon of $N$ steps is:
%
\begin{subequations}
\label{eq:ocp}
\begin{align}
    \underset{\vect{x}_{\mathrm{a},0:N},\,\vect{u}_{\mathrm{a},0:N-1}}
    {\mathrm{minimize}}
        &\quad \sum_{k=0}^{N-1} \ell(\vect{x}_{\mathrm{a},k}, \vect{u}_{\mathrm{a},k})
        + \ell_e(\vect{x}_{\mathrm{a},N})
        \label{eq:ocp_cost}\\
    \mathrm{subject\ to}\quad
        &\vect{x}_{\mathrm{a},k+1} =
        \vect{F}(\vect{x}_{\mathrm{a},k}, \vect{u}_{\mathrm{a},k}),
        \quad k = 0,\ldots,N-1,
        \label{eq:ocp_dyn}\\
        &s_{k+1} = s_k + T_s\,v_{s,k},
        \quad k = 0,\ldots,N-1,
        \label{eq:ocp_arc}\\
        &v_{s,k+1} = v_{s,k} + T_s\,a_{s,k},
        \quad k = 0,\ldots,N-1,
        \label{eq:ocp_progress}\\
        &\vect{x}_{\mathrm{a},0} = \vect{x}_{\mathrm{a},\mathrm{init}},
        \label{eq:ocp_init}\\
        &f_{\min} \le f \le f_{\max},
        \label{eq:ocp_thrust}\\
        &|\tau_i| \le \tau_{\max},
        \quad i = 1,2,3,
        \label{eq:ocp_torque}\\
        &a_{s,\min} \le a_s \le a_{s,\max},
        \label{eq:ocp_speed}\\
        &0 \le s \le s_{\max},
        \qquad
        0 \le v_s \le v_{s,\max},
        \label{eq:ocp_arc_bounds}
\end{align}
\end{subequations}
%
where $\vect{F}(\vect{x}_{\mathrm{a},k},\,\vect{u}_{\mathrm{a},k})$
denotes one step of the augmented dynamics~\eqref{eq:dynamics}
and~\eqref{eq:arc_dyn} integrated by an implicit Runge-Kutta scheme.
The bound constraint~\eqref{eq:ocp_arc_bounds} confines the virtual
reference point to the valid domain of the trajectory interpolant,
preventing extrapolation beyond $s_{\max}$ while ensuring that the
virtual progress speed remains dynamically consistent with the
second-order progress model. 

Although the DQ-MPCC formulation is derived using the nominal
torque-driven quadrotor model, the Software-in-the-Loop validation
described in Section~\ref{sec:mujoco_sil} employs a practical
body-rate-and-thrust command interface consistent with the onboard
flight stack used in the real platform.

\section{Bilevel Optimisation for Weight Tuning}
\label{sec:bilevel}

The stage cost in~\eqref{eq:cost} depends on a vector of tuning weights,
collected as
\begin{equation}
  \boldsymbol{\vartheta} =
  \bigl[\,
    \bm{q}_{c}^\top,\;
    \bm{q}_{\ell}^\top,\;
    \bm{q}_{q}^\top,\;
    \bm{w}_u^\top,\;
    \bm{q}_{\omega}^\top,\;
    Q_s
  \,\bigr]^\top \in \mathbb{R}^{17},
  \label{eq:weight_vector}
\end{equation}
with
$\bm{q}_{c} = [Q_{c}^x, Q_{c}^y, Q_{c}^z]^\top$,
$\bm{q}_{\ell} = [Q_{\ell}^x, Q_{\ell}^y, Q_{\ell}^z]^\top$,
$\bm{q}_{q} = [Q_q^x, Q_q^y, Q_q^z]^\top$,
$\bm{w}_u = [U^f, U^{\tau_1}, U^{\tau_2}, U^{\tau_3}]^\top$,
and
$\bm{q}_{\omega} = [Q_\omega^x, Q_\omega^y, Q_\omega^z]^\top$.
The associated diagonal weighting matrices are defined as
$\bm{Q}_\ell = \mathrm{diag}(\bm{q}_{\ell})$,
$\bm{Q}_c = \mathrm{diag}(\bm{q}_{c})$,
$\bm{Q}_q = \mathrm{diag}(\bm{q}_q)$,
$\bm{Q}_\omega = \mathrm{diag}(\bm{q}_\omega)$,
and
$\bm{W}_u = \mathrm{diag}(\bm{w}_u)$.
The additional scalar $Q_a$ appearing in~\eqref{eq:cost} is kept fixed
throughout the bilevel optimisation in order to isolate the effect of the
main tracking, effort, and progress weights on the closed-loop behaviour.

Since these parameters are mutually coupled and may span different orders of
magnitude, manual tuning alone is difficult to reproduce. In this work the
weight-selection workflow is therefore split into two stages. First, a coarse
manual adjustment is carried out until a stable and practically useful
closed-loop behaviour is obtained, yielding a reference vector
$\boldsymbol{\vartheta}_{\mathrm{ref}}$. Second, bilevel optimisation is used
only as a \emph{local} refinement around this reference point. This preserves
the qualitative behaviour obtained by engineering judgement while avoiding an
overly broad search domain.

During this local refinement, all entries of $\boldsymbol{\vartheta}$ except
$Q_a$ are allowed to vary inside a bounded neighbourhood of the manual
reference. The outer level therefore acts on the local decision vector
\begin{equation}
  \boldsymbol{\psi} =
  \bigl[\,
    \bm{q}_{c}^\top,\;
    \bm{q}_{\ell}^\top,\;
    \bm{q}_{q}^\top,\;
    \bm{w}_u^\top,\;
    \bm{q}_{\omega}^\top,\;
    Q_s
  \,\bigr]^\top \in \mathbb{R}^{17},
  \label{eq:local_weight_vector}
\end{equation}
for the Euclidean MPCC baseline. For the DQ-MPCC controller, the same
local-refinement logic is used, replacing the rotational block $\bm{q}_{q}$
by the corresponding Lie-algebra weight vector $\bm{q}_{\phi}$. In both
cases, the only cost coefficient excluded from the outer decision vector is
$Q_a$.

\subsection{Bilevel Problem Formulation}
\label{sec:bilevel_problem}

Let $\mathcal{V}$ denote a representative set of operating conditions. For each
candidate vector $\boldsymbol{\psi}$ and each $v\in\mathcal{V}$, the
\emph{inner-level} problem consists of solving the OCP in
Section~\ref{sec:ocp_formulation} at each sampling instant and simulating the
resulting closed-loop rollout over the full trajectory. From this rollout, the
outer level extracts contouring and lag RMSE values together with completion,
solver-failure, and path-speed metrics.

The \emph{outer-level} optimisation problem is then written as
\begin{equation}
  \boldsymbol{\psi}^{\star}
  =
  \arg\min_{\boldsymbol{\psi}\in\Theta_{\mathrm{loc}}}
  J_{\mathrm{multi}}(\boldsymbol{\psi}),
  \qquad
  J_{\mathrm{multi}}(\boldsymbol{\psi})
  =
  \frac{1}{|\mathcal{V}|}
  \sum_{v\in\mathcal{V}}
  J(\boldsymbol{\psi},v),
  \label{eq:bilevel_multi}
\end{equation}
where $\Theta_{\mathrm{loc}}$ denotes a local neighbourhood around the manual
reference controller $\boldsymbol{\vartheta}_{\mathrm{ref}}$.

For both controllers, each sub-objective is defined as
\begin{equation}
\begin{aligned}
J(\boldsymbol{\psi},v)
&=
W_{c}\,\frac{\mathrm{RMSE}_{c}}{\mathrm{RMSE}_{c}^{\mathrm{ref}}(v)}
+ W_{\ell}\,\frac{\mathrm{RMSE}_{\ell}}{\mathrm{RMSE}_{\ell}^{\mathrm{ref}}(v)}
\\[2pt]
&\quad
- W_{v}\,\frac{v_{\mathrm{eff}}}{v_{\mathrm{eff}}^{\mathrm{ref}}(v)}
+ W_{\mathrm{inc}}\,[1-c]^+
+ W_{\mathrm{fail}}\,\eta_{\mathrm{fail}},
\end{aligned}
\label{eq:sub_objective}
\end{equation}
where $[\cdot]^+=\max(\cdot,0)$,
$c=\theta_K/s_{\max}$ is the path-completion ratio,
$\eta_{\mathrm{fail}}$ is the mean solver-failure ratio, and
\begin{equation}
  \bar{\rho}
  =
  \frac{1}{K}\sum_{k=0}^{K-1}\frac{v_{\mathrm{path},k}}{v_{s,k}+\varepsilon},
  \qquad
  v_{\mathrm{eff}} = \bar{v}_s\,\bar{\rho},
  \label{eq:effective_speed}
\end{equation}
with $\bar{v}_s$ the mean virtual progress speed and $v_{\mathrm{path},k}$ the
realised path-relative speed at step $k$. The superscript ``ref'' denotes the
metrics obtained with the manually tuned controller
$\boldsymbol{\vartheta}_{\mathrm{ref}}$ at the same operating condition $v$.

These upper-level costs are aligned with the practical tuning goal of the
paper: reduce contouring and lag while preserving a physically meaningful
relation between virtual progress and actual spatial motion. The term based on
$v_{\mathrm{eff}}$ rewards fast traversal only when the virtual progress
remains consistent with the realised path speed, thereby discouraging solutions
that improve the arc-length state artificially without improving the true
path-following behaviour.

The outer-level optimisation is solved with Bayesian optimisation, while the
inner level is the full closed-loop controller evaluation.

\section{Experimental Setup}
\label{sec:impl}

%\subsection{Platform and Hardware}

Numerical simulations are carried out on a desktop workstation equipped with an
AMD Ryzen 7 7000 series processor running Ubuntu 22.04, while the experimental
validation is conducted on a custom quadrotor platform with an NVIDIA Jetson Orin
NX 16\,GB as onboard computer.

The high-level control loop runs at \SI{100}{\hertz}, and the solver
consistently executes below this timing bound for a prediction horizon of
\SI{0.3}{\second}. The second-order progress model is handled inside the
same OCP, so no additional smoothing layer is required outside the solver.
A more detailed analysis of the effect of the prediction horizon and the
number of discretization steps \(N\) is left for future work.
This setup enables evaluation of the full closed-loop software stack under
flight-representative computational conditions while preserving consistency with
the command pathway used in the real platform. The corresponding simulator
approximation is described next.

The OCP~\eqref{eq:ocp} is solved within the ACADOS framework~\cite{verschueren2021acados},
employing a Sequential Quadratic Programming Real-Time Iteration (SQP-RTI) scheme that
executes a single QP per control cycle~\cite{diehl2005real}. Quadratic subproblems are handled by HPIPM with
full condensing, while system dynamics are integrated via an Implicit Runge-Kutta (IRK)
scheme. The Hessian is approximated via Gauss-Newton with CONVEXIFY regularization, and
stage costs are expressed as CasADi symbolic externals. The controller operates at a
sample time of $T_s = \SI{10}{\milli\second}$ over a prediction horizon of
$T_N = \SI{0.3}{\second}$, discretised into $N = 30$ shooting nodes.

A critical implementation aspect is the representation of the reference
trajectory $\bm{\gamma}(s)$ as a CasADi symbolic function compatible with
automatic differentiation in ACADOS. The map is constructed from
$N_w = 30$ arc-length waypoints over $s \in [0, s_{\max}]$ using
piecewise-linear interpolation of the dual-quaternion and velocity
components. Quaternion continuity is preserved by enforcing hemisphere
consistency, i.e., flipping $q_{r,i}^d \leftarrow -q_{r,i}^d$ whenever
$q_{r,i-1}^d \cdot q_{r,i}^d < 0$, while the heading angle is unwrapped
before interpolation.

%\subsection{Weight-Tuning Setup}
%\label{sec:weight_tuning_setup}

For the experiments in this work, the bilevel tuning procedure is instantiated
using the representative velocity set
\[
\mathcal{V}=\{5,\;10,\;16\}\ \text{m/s}.
\]
Each tuning rollout uses the same headless MiL configuration: maximum simulated
duration $T_{\mathrm{final}}=\SI{30}{\second}$, trajectory-definition window
$T_{\mathrm{traj}}=\SI{30}{\second}$, control frequency
$f_s=\SI{100}{\hertz}$ ($T_s=\SI{10}{\milli\second}$), and prediction horizon
$T_N=\SI{0.3}{\second}$ discretised into $N=30$ shooting nodes with
$N_w=30$ arc-length interpolation waypoints. Under this setup, one rollout
contains approximately $2971$ closed-loop simulation steps, and each Optuna
trial evaluates one full rollout for every operating point in $\mathcal{V}$.

After coarse manual tuning, the local search domain is defined around the
reference controller $\boldsymbol{\vartheta}_{\mathrm{ref}}$ using
hierarchical multiplicative windows. The tracking and orientation blocks
$\bm{q}_{c}$, $\bm{q}_{\ell}$, and $\bm{q}_{q}$ or $\bm{q}_{\phi}$ are searched
within $[0.75,\;1.35]\times(\cdot)^{\mathrm{ref}}$, the progress weight is
searched within $[0.75,\;1.35]\times Q_{s}^{\mathrm{ref}}$, the effort block
$\bm{w}_u$ is searched within $[0.92,\;1.08]\times(\bm{w}_u)^{\mathrm{ref}}$,
and the angular-rate block $\bm{q}_{\omega}$ is searched within
$[0.90,\;1.10]\times(\bm{q}_{\omega})^{\mathrm{ref}}$. The progress-acceleration
weight $Q_a$ remains fixed at its manually tuned value. This choice avoids
confounding the comparison of tracking-oriented cost weights with the internal
smoothing imposed on the second-order virtual-progress dynamics, so the bilevel
refinement focuses on the blocks that directly shape path-following behaviour.
The resulting local search space is reported in Table~\ref{tab:search_space}.

The same upper-level coefficients are used for both controllers, so that the
outer bilevel comparison remains methodologically consistent across DQ-MPCC and
the Euclidean MPCC baseline. In the present study, the local objective
in~\eqref{eq:sub_objective} is instantiated with
\(
W_{c}=1.5,\;
W_{\ell}=0.35,\;
W_{v}=0.80,\;
W_{\mathrm{inc}}=1000,\;
W_{\mathrm{fail}}=500.
\)

The outer optimisation is performed with the Tree-structured Parzen Estimator
(TPE) algorithm implemented in Optuna~\cite{akiba2019optuna}, using 10 uniform
warm-up trials and random seed 42 for reproducibility. Since each trial requires
one full closed-loop simulation for each operating point in $\mathcal{V}$, the
computational cost of each outer iteration is approximately
$|\mathcal{V}|$ times that of a single simulation run. In the current study,
the final local refinement stage is carried out with 100 TPE trials per
controller, which provides sufficient exploration of the bounded neighbourhood
around the manually tuned operating point without resorting to a full
high-dimensional global search.

\begin{table}[H]
\centering
\caption{Local search window $\Theta_{\mathrm{loc}}$ used in the bilevel
refinement around the manually tuned reference controller
($i\in\{x,y,z\}$).}
\label{tab:search_space}
\renewcommand{\arraystretch}{1.1}
\setlength{\tabcolsep}{6pt}
\begin{tabular}{l c}
\toprule
\textbf{Parameter} & \textbf{Local range} \\
\midrule
$Q_{c,i}$        & $[0.75,\;1.35]\times Q_{c,i}^{\mathrm{ref}}$ \\
$Q_{\ell,i}$     & $[0.75,\;1.35]\times Q_{\ell,i}^{\mathrm{ref}}$ \\
$Q_{q,i}$ or $Q_{\phi,i}$ & $[0.75,\;1.35]\times Q_{q,i}^{\mathrm{ref}}$ or $Q_{\phi,i}^{\mathrm{ref}}$ \\
$Q_s$            & $[0.75,\;1.35]\times Q_{s}^{\mathrm{ref}}$ \\
$U^f,\;U^{\tau_j}$ & $[0.92,\;1.08]\times (U^f,\;U^{\tau_j})^{\mathrm{ref}}$ \\
$Q_{\omega,i}$   & $[0.90,\;1.10]\times Q_{\omega,i}^{\mathrm{ref}}$ \\
$Q_a$            & fixed at manual reference \\
\bottomrule
\end{tabular}
\end{table}




%\subsection{Reference Trajectory}

The reference path is a three-dimensional Lissajous curve, defined
in closed form as
\begin{equation}
    \vect{p}^d(t)
    = \begin{bmatrix}
        A_x \sin(\Omega_x t + \delta_x)\\
        A_y \sin(\Omega_y t)\\
        A_z \cos(\Omega_z t) + z_0
      \end{bmatrix},
    \label{eq:lissajous}
\end{equation}
where $A_x, A_y, A_z$ are the trajectory amplitudes,
$\Omega_x, \Omega_y, \Omega_z$ the path frequencies,
$\delta_x$ a phase offset that breaks planar symmetry,
and $z_0$ the nominal flight altitude. The curve is periodic
with period $T_L$ and combines low- and high-curvature segments
that exercise the lag--contouring decomposition across a
representative range of flight conditions. The time-parametrized
path $\vect{p}^d(t)$ is evaluated analytically, so derivatives
$\dot{\vect{p}}^d$ and $\ddot{\vect{p}}^d$ are available in
closed form without numerical differentiation. The specific
parameter values used in all experiments are reported in
Table~\ref{tab:lissajous_params}.

All simulations are conducted on the three-dimensional Lissajous
path defined in Section~\ref{sec:ref_traj}, equation~\eqref{eq:lissajous},
with the parameters reported in Table~\ref{tab:lissajous_params}.
The resulting curve spans a total arc length of approximately
\SI{80}{\meter} over a period of $T_L = \SI{15}{\second}$,
combining low-curvature straight segments with tight turns that
stress the lag--contouring decomposition across both controllers.

\begin{table}[h]
\centering
\caption{Lissajous Trajectory Parameters}
\label{tab:lissajous_params}
\begin{tabular}{lc}
\toprule
\textbf{Parameter} & \textbf{Value} \\
\midrule
Amplitudes $(A_x,\,A_y,\,A_z)$                 & $(10.0,\;8.0,\;3.0)$ \si{\meter} \\
Frequencies $(\Omega_x,\,\Omega_y,\,\Omega_z)$ & $(0.42,\;0.84,\;1.26)$ \si{\radian\per\second} \\
Phase offset $\delta_x$                         & $\pi/2$ \si{\radian} \\
Nominal altitude $z_0$                          & $5.0$ \si{\meter} \\
Period $T_L$                                    & $15.0$ \si{\second} \\
Arc length $s_{\max}$                           & ${\approx}\,80$ \si{\meter} \\
\bottomrule
\end{tabular}
\end{table}




% ═══════════════════════ V. RESULTS ════════════════════════
\section{Results} 
\label{sec:results}



\subsection{MiL Simulations}

\subsubsection{Local Tuning Report}

The local refinement reported here follows the common bilevel formulation and
experimental setup defined in Section~\ref{sec:bilevel_problem} and
Section~\ref{sec:impl}. Table~\ref{tab:tuning_compact_combined} reports the
manual and refined gains selected by the local search, while
Table~\ref{tab:tuning_manual_vs_refined} summarises the corresponding
closed-loop effect in the final MiL baseline runs.

\begin{table}[H]
\centering
\caption{Local tuning summary for both controllers.}
\label{tab:tuning_compact_combined}
\scriptsize
\renewcommand{\arraystretch}{1.05}
\setlength{\tabcolsep}{2pt}
\resizebox{\columnwidth}{!}{%
\begin{tabular}{lccc}
\toprule
\textbf{Weight} & \textbf{Manual} & \textbf{DQ-MPCC} & \textbf{MPCC} \\
\midrule
$\bm{q}_{\phi}/\bm{q}_{q}$ 
& $[0.65, 0.65, 1.9]$ 
& $[0.53, 0.69, 1.7]$ 
& $[0.62, 0.72, 1.6]$ \\
$\bm{q}_{c}$ 
& $[40.0, 40.0, 40.0]$ 
& $[44.4, 37.9, 32.6]$ 
& $[40.6, 35.9, 33.7]$ \\
$\bm{q}_{\ell}$ 
& $[3.0, 3.0, 3.0]$ 
& $[4.0, 2.6, 3.6]$ 
& $[4.0, 4.0, 2.7]$ \\
$\bm{w}_{u}$ 
& $[0.050, 70.0, 70.0, 70.0]$ 
& $[0.049, 68.9, 69.6, 66.6]$ 
& $[0.049, 70.9, 67.6, 72.2]$ \\
$\bm{q}_{\omega}$ 
& $[0.020, 0.020, 0.10]$ 
& $[0.021, 0.021, 0.100]$ 
& $[0.020, 0.021, 0.095]$ \\
$Q_s$ 
& $11.0$ 
& $8.4$ 
& $8.3$ \\
\midrule
$J^\star$ & -- & 0.917 & 0.863 \\
Trial & -- & \#72 & \#94 \\
\bottomrule
\end{tabular}%
}
\end{table}

\begin{table}[H]
\centering
\caption{Closed-loop comparison between manual and refined gains used in the final baseline runs.}
\label{tab:tuning_manual_vs_refined}
\footnotesize
\renewcommand{\arraystretch}{1.12}
\setlength{\tabcolsep}{3.5pt}
\begin{tabular}{lcccccc}
\toprule
\textbf{Controller} & \textbf{Variant} & $\mathbf{RMSE}_c$ & $\mathbf{RMSE}_\ell$ & $\mathbf{v}_{\mathrm{eff}}$ & \textbf{Path} & \textbf{Fail} \\
\midrule
DQ-MPCC & Manual  & 0.346 & 0.507 & 1.092 & 100.00\% & 0.00\% \\
DQ-MPCC & Refined & 0.291 & 0.335 & 1.216 & 100.00\% & 0.00\% \\
\cmidrule(lr){2-7}
DQ-MPCC & $\Delta$ (\%) & 16.1 & 33.9 & 11.3 & -- & -- \\
\midrule
MPCC    & Manual  & 0.445 & 0.974 & 5.134 & 99.98\% & 0.00\% \\
MPCC    & Refined & 0.337 & 0.476 & 3.979 & 99.98\% & 0.00\% \\
\cmidrule(lr){2-7}
MPCC    & $\Delta$ (\%) & 24.4 & 51.1 & -22.5 & -- & -- \\
\bottomrule
\end{tabular}
\end{table}

Overall, both controllers benefited from the local refinement. While the resulting closed-loop effects differ
between controllers, this stage is important to ensure that the subsequent comparison is conducted on a fairer
and more balanced basis, since both controllers are evaluated after the same consistent local tuning procedure.

\subsubsection{Velocity Sweep: Contouring and Lag Error Analysis}
\label{sec:velocity_sweep_analysis}

To rigorously evaluate the tracking quality of both formulations across
operating regimes, we decompose the position tracking error into its
orthogonal components: the \emph{contouring error} $e_c$ (perpendicular
distance to the reference path) and the \emph{lag error} $e_\ell$
(signed projection along the path tangent).  By the Pythagorean
relationship these satisfy
\begin{equation}
    \mathrm{RMSE}_{e_p}^{2} = \mathrm{RMSE}_{e_c}^{2} + \mathrm{RMSE}_{e_\ell}^{2},
    \label{eq:decomposition}
\end{equation}
which was verified numerically for every run in our Monte~Carlo campaign.
Both components are computed in the inertial frame to ensure a fair,
frame-neutral comparison between the DQ-MPCC (which internally operates
in the body frame via $\mathrm{Log}\bigl(\mathrm{SE}(3)\bigr)$) and
the Baseline MPCC (inertial-frame errors).

The experiment considers eight maximum progress velocities,
\[
v_\theta^{\max} \in \{4,\,6,\,8,\,10,\,12,\,15,\,16,\,18\}\ \text{m/s},
\]
with 10 Monte Carlo runs per velocity and per controller
($\sigma_p = 0.05$~m, $\sigma_q = 0.05$~rad), for a total of
160 simulations along a Lissajous path of arc length
$s_{\max} = 80$~m.

%% ─────────────────────────────────────────────────────────────────────
%%  TABLE I — Full comparative metrics
%% ─────────────────────────────────────────────────────────────────────

\begin{table*}[t]
\centering
\caption{Comparative performance across progress velocities.
$\Delta$\% is computed as $(x_{\text{DQ}} - x_{\text{Base}})/x_{\text{Base}} \times 100$;
negative values favour DQ-MPCC.  Bold marks the better value per velocity.}
\label{tab:velocity_sweep_full}
\renewcommand{\arraystretch}{1.1}
\setlength{\tabcolsep}{3pt}
\resizebox{\textwidth}{!}{%
\begin{tabular}{r|ccc|ccc|ccc|ccc|ccc}
\toprule
& \multicolumn{3}{c|}{Contouring $\mathrm{RMSE}_{e_c}$ [m]}
& \multicolumn{3}{c|}{Lag $\mathrm{RMSE}_{e_\ell}$ [m]}
& \multicolumn{3}{c|}{Position $\mathrm{RMSE}_{e_p}$ [m]}
& \multicolumn{3}{c|}{Orientation $\mathrm{RMSE}_{e_\psi}$ [rad]}
& \multicolumn{3}{c}{Lap time $t_\mathrm{lap}$ [s]} \\
$v_\theta^{\max}$
& DQ & Base & $\Delta$\%
& DQ & Base & $\Delta$\%
& DQ & Base & $\Delta$\%
& DQ & Base & $\Delta$\%
& DQ & Base & $\Delta$\% \\
\midrule
     4 & \textbf{0.2863} & 0.3287 & -12.9\% & 0.2698 & \textbf{0.2168} & +24.4\% & \textbf{0.3934} & 0.3938 & -0.1\% & \textbf{0.3808} & 0.4972 & -23.4\% & 22.29 & 24.03 & -7.3\% \\
 6 & \textbf{0.3294} & 0.3850 & -14.4\% & 0.3332 & \textbf{0.2719} & +22.5\% & \textbf{0.4685} & 0.4713 & -0.6\% & \textbf{0.5338} & 0.6767 & -21.1\% & 15.60 & 16.47 & -5.3\% \\
 8 & \textbf{0.3869} & 0.4549 & -14.9\% & 0.4059 & \textbf{0.3083} & +31.6\% & 0.5608 & \textbf{0.5498} & +2.0\% & \textbf{0.6561} & 0.8330 & -21.2\% & 12.24 & 12.71 & -3.8\% \\
10 & \textbf{0.4340} & 0.4521 & -4.0\% & 0.4705 & \textbf{0.3495} & +34.6\% & 0.6401 & \textbf{0.5715} & +12.0\% & \textbf{0.7520} & 0.8956 & -16.0\% & 10.18 & 10.10 & +0.8\% \\
12 & \textbf{0.4793} & 0.4800 & -0.1\% & 0.5335 & \textbf{0.3904} & +36.6\% & 0.7172 & \textbf{0.6188} & +15.9\% & \textbf{0.8239} & 0.9995 & -17.6\% & 8.74 & 8.81 & -0.7\% \\
15 & 0.5312 & \textbf{0.5118} & +3.8\% & 0.6404 & \textbf{0.4262} & +50.3\% & 0.8321 & \textbf{0.6660} & +24.9\% & \textbf{0.9225} & 1.0996 & -16.1\% & 7.37 & 7.42 & -0.6\% \\
16 & 0.5482 & \textbf{0.5479} & +0.1\% & 0.6764 & \textbf{0.4391} & +54.0\% & 0.8707 & \textbf{0.7021} & +24.0\% & \textbf{0.9490} & 1.1123 & -14.7\% & 7.05 & 7.05 & -0.0\% \\
18 & \textbf{0.5828} & 0.5929 & -1.7\% & 0.7504 & \textbf{0.4738} & +58.4\% & 0.9502 & \textbf{0.7590} & +25.2\% & \textbf{0.9902} & 1.1462 & -13.6\% & 6.51 & 6.49 & +0.4\% \\
\midrule
    \textit{Avg.} & 0.4473 & 0.4692 & -4.7\% & 0.5100 & 0.3595 & +41.9\% & 0.6791 & 0.5915 & +14.8\% & 0.7510 & 0.9075 & -17.2\% & 11.25 & 11.63 & -3.3\% \\
\bottomrule
\end{tabular}}
\end{table*}

%% ─────────────────────────────────────────────────────────────────────
%%  TABLE II — Orthogonal decomposition
%% ─────────────────────────────────────────────────────────────────────

\begin{table}[H]
\centering
\caption{Orthogonal decomposition of the position error into
contouring (\%$e_c^2$) and lag (\%$e_\ell^2$) components.
The ratio $e_\ell/e_c$ indicates the dominant error type:
$>1$ implies lag-dominated tracking, $<1$ implies contouring-dominated.}
\label{tab:decomposition}
\renewcommand{\arraystretch}{1.1}
\begin{tabular}{r|ccc|ccc}
\toprule
& \multicolumn{3}{c|}{DQ-MPCC}
& \multicolumn{3}{c}{Baseline MPCC} \\
$v_\theta^{\max}$ & \%$e_c^2$ & \%$e_\ell^2$ & $e_\ell/e_c$
                  & \%$e_c^2$ & \%$e_\ell^2$ & $e_\ell/e_c$ \\
\midrule
     4 & 53.0 & 47.0 & 0.942 & 69.7 & 30.3 & 0.660 \\
 6 & 49.4 & 50.6 & 1.011 & 66.7 & 33.3 & 0.706 \\
 8 & 47.6 & 52.4 & 1.049 & 68.6 & 31.4 & 0.678 \\
10 & 46.0 & 54.0 & 1.084 & 62.6 & 37.4 & 0.773 \\
12 & 44.7 & 55.3 & 1.113 & 60.2 & 39.8 & 0.813 \\
15 & 40.8 & 59.2 & 1.206 & 59.1 & 40.9 & 0.833 \\
16 & 39.7 & 60.3 & 1.234 & 60.9 & 39.1 & 0.801 \\
18 & 37.6 & 62.4 & 1.287 & 61.0 & 39.0 & 0.799 \\
\bottomrule
\end{tabular}
\end{table}

%% ─────────────────────────────────────────────────────────────────────
%%  FIGURE — Pareto frontiers (2×2)
%% ─────────────────────────────────────────────────────────────────────

\begin{figure}[H]
\centering
\includegraphics[width=\columnwidth]{Fig/fig_pareto_quad.pdf}
\caption{Speed--accuracy Pareto frontiers for the DQ-MPCC and
Baseline MPCC across 8 progress velocities.
Each point is the median over 10 Monte~Carlo runs;
error bars denote the interquartile range (P25--P75).
(a)~Contouring RMSE, (b)~Lag RMSE, (c)~Position RMSE,
(d)~Orientation RMSE, all vs.\ lap time.
Annotations indicate $v_\theta^{\max}$ [m/s].}
\label{fig:pareto_quad}
\end{figure}

%% ─────────────────────────────────────────────────────────────────────
%%  FIGURE ANALYSIS
%% ─────────────────────────────────────────────────────────────────────

Fig.~\ref{fig:pareto_quad} presents the speed--accuracy Pareto frontiers
for both controllers.
%
In panel~(a), the DQ-MPCC (blue) lies below the Baseline (red) for $v_\theta^{\max} < 15$~m/s, after which the curves nearly merge (gap $<0.010$~m at $18$~m/s),
confirming that the coupled $\mathrm{SE}(3)$ formulation yields tighter
geometric adherence to the path.
%
Panel~(b) reveals the opposite: the Baseline's lag curve is consistently
lower, with an absolute gap that widens from
$0.053$~m at $4$~m/s
to $0.277$~m at $18$~m/s,
reflecting the decoupled formulation's ability to independently minimise
along-path delay.
%
Panel~(c) combines both components via~\eqref{eq:decomposition}:
the curves cross near $v_\theta^{\max} \approx 8$~m/s, with the lag-dominated DQ-MPCC yielding higher total position error
at high speeds.
%
Finally, panel~(d) shows that the DQ-MPCC maintains consistently lower
orientation error across the entire velocity range
($13.6$--$23.4\%$),
a direct benefit of the pose-level coupling in the cost function.


\begin{figure}[H]
\centering
\includegraphics[width=\columnwidth]{Fig/fig_3d_velocity.pdf}
\caption{Three-dimensional closed-loop trajectories for both controllers
across multiple progress-velocity limits. Grey curves correspond to
$v_{\theta}^{\max} \in \{4,6,8,12,14,16,18,20\}$~m/s, while the red curve
highlights $v_{\theta}^{\max}=10$~m/s as the nominal operating point. The
dashed black curve is the arc-length-parameterised reference path
($s_{\max}=80$~m). The DQ-MPCC maintains tighter spatial adherence in
high-curvature regions, whereas the baseline tends to achieve smaller lag
at the expense of larger orientation and contouring deviations.}
\label{fig:3d_beahaibiur}
\end{figure}

%% ─────────────────────────────────────────────────────────────────────
%%  ANALYSIS  (all numbers auto-generated)
%% ─────────────────────────────────────────────────────────────────────

\subsubsection{Contouring vs.\ Lag Decomposition}

Table~\ref{tab:velocity_sweep_full} shows that the DQ-MPCC achieves lower
contouring error for $v_\theta^{\max} < 15$~m/s
(avg.\ $-4.7\%$, best $-14.9\%$ at $8$~m/s),
whereas the Baseline MPCC dominates in lag error
(avg.\ $+41.9\%$, up to $+58.4\%$ at $18$~m/s).
The decomposition in Table~\ref{tab:decomposition} explains this asymmetry:
the Baseline is \emph{contouring-dominated}
($59$--$70\%$ of $\|e_p\|^2$ comes from $e_c$,
$e_\ell/e_c < 1$ at all speeds),
while the DQ-MPCC transitions from near parity
($53$/$\ 47$\%
at $4$~m/s) to \emph{lag-dominated}
($38$/$\ 62$\%,
$e_\ell/e_c = 1.29$ at $18$~m/s).
%
The root cause is structural: the DQ-MPCC's $\mathrm{SE}(3)$ logarithmic
error map $\boldsymbol{\rho} = \mathbf{J}_l(\boldsymbol{\varphi})^{-1}
\mathbf{R}_d^\top \Delta\mathbf{p}$ couples position and orientation,
preventing the solver from independently minimising the tangential projection.
The Baseline, by contrast, penalises $e_c$ and $e_\ell$ as decoupled cost
terms, enabling aggressive lag reduction at the expense of orientation fidelity.
These results therefore do not support a blanket claim of superiority of one
controller over the other; rather, they show that the proposed DQ-MPCC
systematically shifts the closed-loop trade-off toward better contouring and
attitude consistency, while the Euclidean baseline remains more effective at
reducing tangential lag.

% ═══════════════════════════════════════════════════════════════════
%  AUTO-GENERATED — do not edit manually
%  Source: latex/generate_montecarlo.py
% ═══════════════════════════════════════════════════════════════════


\subsection{Experiment~3: Monte Carlo Robustness Evaluation}
\label{sec:experiment3}

To evaluate robustness under uncertain initial conditions, we perform
a Monte Carlo study with $N = 100$ trials at a fixed progress-velocity
limit $v_{\theta}^{\max} = 10$~m/s.  Each trial samples
a random initial pose: position is perturbed as
$\mathbf{p}_0 \sim \mathcal{N}(\bar{\mathbf{p}}_0,\,
\sigma_p^2 \mathbf{I}_3)$ with $\sigma_p = 2.0$~m,
and orientation as a random rotation with
$\|\mathrm{Log}(\delta q)\| \sim U(0,\, 0.5)$~rad
($\approx 29^\circ$).
The path length is $s_{\max} = 80$~m with a time budget
of $T = 40$~s.
A trial is deemed \emph{successful} if the progress variable reaches
$\theta \geq 95\%$ of $s_{\max}$.

\begin{table}[H]
  \centering
  \caption{Monte Carlo robustness comparison ($v = 10$~m/s, $N = 100$).}
  \label{tab:montecarlo}
  \setlength{\tabcolsep}{4pt}
  \begin{tabular}{l c c c}
    \toprule
    \textbf{Metric} & \textbf{DQ-MPCC} & \textbf{Baseline} & $\boldsymbol{\Delta}$ \\
    \midrule
    Conv. rate [\%] & 99.0 & 100.0 & -1.0\% \\
    RMSE $e_c$ & $0.8243 \pm 0.5405$ & $0.7632 \pm 0.6192$ & +8.0\% \\
    RMSE $e_l$ & $0.6425 \pm 0.3936$ & $0.4827 \pm 0.2650$ & +33.1\% \\
    RMSE $e_{\mathrm{pos}}$ & $1.0824 \pm 0.6066$ & $0.9390 \pm 0.6224$ & +15.3\% \\
    RMSE $e_{\mathrm{ori}}$ & $0.9118 \pm 0.1147$ & $1.0132 \pm 0.0419$ & -10.0\% \\
    \midrule
    Solver time [ms] & $3.81 \pm 0.11$ & $2.04 \pm 0.09$ & +86.7\% \\
    Ctrl. effort & $33518.1 \pm 75505.4$ & $35187.1 \pm 30591.3$ & -4.7\% \\
    Lap time [s] & $10.57 \pm 0.29$ & $10.39 \pm 0.30$ & +1.7\% \\
    \bottomrule
  \end{tabular}
\end{table}

\paragraph{Convergence}
The baseline achieves a convergence rate of 100.0\%, compared to 99.0\% for DQ-MPCC.
Both controllers achieve near-perfect convergence despite large
initial-pose perturbations ($\sigma_p = 2.0$~m,
$\sigma_q = 0.5$~rad), confirming reliable operation
under uncertainty.

\begin{figure}[!t]
  \centering
  \includegraphics[width=\columnwidth]{Fig/fig_mc_robustness.pdf}
  \caption{Distribution of per-run RMSE across 100 Monte Carlo trials
           with random initial poses ($\sigma_p = 2.0$~m,
           $\sigma_q = 0.5$~rad) at $v_{\theta}^{\max} = 10$~m/s.
           Violin contours show probability density; individual points
           overlay each trial.  Median values are annotated.}
  \label{fig:mc_robustness}
\end{figure}

\paragraph{Tracking accuracy and Pareto consistency}
Fig.~\ref{fig:mc_robustness} displays the error distributions
for all four RMSE metrics.
The orientation subplot confirms that DQ-MPCC consistently reduces
$e_{\mathrm{ori}}$: its median
($0.8708)$ lies well below the baseline median
($1.0046)$, with DQ-MPCC winning in
85\,/\,99 individual trials
($\Delta\bar{e}_{\mathrm{ori}} = -10.0\%$).
Crucially, the baseline's $e_{\mathrm{ori}}$ distribution
is tightly concentrated ($\sigma = 0.0419$),
while DQ-MPCC's wider spread
($\sigma = 0.1147$) reflects a
pose-dependent adaptation enabled by the Lie-algebraic coupling
between translation and rotation.
%
This trade-off mirrors the Pareto front in Experiment~2
(Fig.~\ref{fig:pareto_quad}): at $v = 10$~m/s with nominal
initial conditions, DQ-MPCC achieves $e_c = 0.434$
($-4.0\%$) and $e_{\mathrm{ori}} = 0.752$
($-16.0\%$) relative to the baseline.
Under large pose perturbations ($\sigma_p = 2.0$~m),
the transient convergence phase inflates the contouring error
to $\bar{e}_c = 0.8243$ (+8.0\%),
yet the orientation advantage is preserved (-10.0\%).
%
Notably, the lag error $e_l$ shows a consistent trend across
both experiments (Exp.~2: $+34.6\%$, Exp.~3: +33.1\%),
indicating that DQ-MPCC systematically trades
progress speed for orientation alignment. This trade-off is desirable
for applications in which rigid-body pose consistency and attitude
quality are more critical than purely tangential progress tracking.

\paragraph{Computational cost}
The mean solver time per MPC step is
$3.81$~ms for DQ-MPCC and
$2.04$~ms for the baseline (+86.7\%).
Both remain well within the $10$~ms real-time budget
($f = 100$~Hz).

\paragraph{Control effort}
The accumulated control-rate effort
$\sum\|\Delta \mathbf{u}\|^2$ averages
$33518.1$ for DQ-MPCC and
$35187.1$ for the baseline (-4.7\%),
indicating smoother
actuation for the dual-quaternion formulation.



\subsection{Simulation Experiment}

\label{sec:mujoco_sil}
Prior to flight tests, both controllers are validated in a MuJoCo-based
Software-in-the-Loop (SiL) setup executed on the Jetson Orin NX. Because the
onboard flight stack accepts body-rate and thrust commands rather than direct
torque inputs, the simulator is configured to reproduce the same control
interface.

Unlike the nominal rigid-body model introduced in
Section~\ref{sec:dynamics}, where the rotational dynamics are expressed
in terms of applied control torques, the MuJoCo simulator employs a
closed-loop approximation of the onboard angular-rate controller.
Specifically, the rotational subsystem is represented by the first-order
response
\begin{equation}
    \dot{\vect{\omega}}
    =
    -\mat{K}_{\omega}
    \left(
        \vect{\omega} - \vect{\omega}_{\mathrm{cmd}}
    \right),
    \label{eq:sil_rate_model}
\end{equation}
where $\vect{\omega} \in \mathbb{R}^{3}$ is the measured body angular
velocity, $\vect{\omega}_{\mathrm{cmd}} \in \mathbb{R}^{3}$ is the
commanded body-rate vector, and
$\mat{K}_{\omega} = \mathrm{diag}(k_x,k_y,k_z)$ is a diagonal gain
matrix that approximates the effective bandwidth of the inner rate loop.

Accordingly, the simulator input is defined as
\begin{equation}
    \vect{u}_{\mathrm{sil}}
    =
    \begin{bmatrix}
        f &
        \omega_{x,\mathrm{cmd}} &
        \omega_{y,\mathrm{cmd}} &
        \omega_{z,\mathrm{cmd}}
    \end{bmatrix}^{\!\top},
    \label{eq:sil_input}
\end{equation}
which matches the practical command interface used by the flight stack.

This SiL approximation is introduced only for implementation-level
validation and does not replace the nominal control-oriented model used
in the DQ-MPCC derivation. Rather, it serves as a closed-loop surrogate
of the onboard rate controller, enabling a more realistic assessment of
timing, numerical behaviour, and command consistency prior to real-world
flight tests.

\subsection{Real Experiment}
% ═══════════════════════ REFERENCES ════════════════════════
\bibliographystyle{IEEEtran}

\begin{thebibliography}{10}

\bibitem{liniger2015optimization}
A.~Liniger, A.~Domahidi, and M.~Morari, ``Optimization-based autonomous racing of 1:43 scale RC cars,'' \emph{Optimal Control Applications and Methods}, vol.~36, no.~5, pp.~628--647, 2015.

\bibitem{romero2022model}
A.~Romero, S.~Sun, P.~Foehn, and D.~Scaramuzza, ``Model predictive contouring control for time-optimal quadrotor flight,'' \emph{IEEE Transactions on Robotics}, vol.~38, no.~6, pp.~3340--3356, 2022.

\bibitem{kabzan2019amz}
J.~Kabzan \emph{et~al.}, ``AMZ Driverless: The full autonomous racing system,'' \emph{Journal of Field Robotics}, vol.~37, no.~7, pp.~1267--1294, 2020.

\bibitem{wang2012}
X.~Wang and C.~Yu, ``Unit dual quaternion-based feedback linearization tracking problem for attitude and position dynamics,'' \emph{Systems \& Control Letters}, vol.~61, no.~12, pp.~1261--1269, 2012.

\bibitem{abaunza2017}
H.~Abaunza, L.~Montes, and C.~Ferreira, ``Dual quaternion modeling and control of a quad-rotor aerial manipulator,'' \emph{IEEE Latin America Transactions}, vol.~15, no.~2, pp.~220--228, 2017.

\bibitem{kenwright2012dual}
B.~Kenwright, ``A beginners guide to dual-quaternions,'' in \emph{Proc. 20th International Conference in Central Europe on Computer Graphics, Visualization and Computer Vision}, 2012, pp.~1--13.

\bibitem{johansen2019review}
T.~A.~Johansen, T.~I.~Fossen, S.~M.~A.~Shah, and A.~S.~de~Oliveira, ``Dual quaternion control: A review of recent results within motion control,'' \emph{Nonlinear Studies}, vol.~26, no.~2, pp.~265--281, 2019.

\bibitem{daniilidis1999}
K.~Daniilidis, ``Hand-eye calibration using dual quaternions,'' \emph{The International Journal of Robotics Research}, vol.~18, no.~3, pp.~286--298, 1999.

\bibitem{selig2005}
J.~M.~Selig, \emph{Geometric Fundamentals of Robotics}, 2nd~ed.\hskip 1em plus 0.5em minus 0.4em\relax Springer, 2005.

\bibitem{sola2018micro}
J.~Sol\`{a}, J.~Deray, and D.~Atchuthan, ``A micro Lie theory for state estimation in robotics,'' \emph{arXiv preprint arXiv:1812.01537}, 2018.

\bibitem{barfoot2017}
T.~D.~Barfoot, \emph{State Estimation for Robotics}.\hskip 1em plus 0.5em minus 0.4em\relax Cambridge University Press, 2017.

\bibitem{verschueren2021acados}
R.~Verschueren \emph{et~al.}, ``acados---A modular open-source framework for fast embedded optimal control,'' \emph{Mathematical Programming Computation}, vol.~14, pp.~147--183, 2022.

\bibitem{murray1994mathematical}
R.~M.~Murray, Z.~Li, and S.~S.~Sastry, \emph{A Mathematical Introduction to Robotic Manipulation}.\hskip 1em plus 0.5em minus 0.4em\relax CRC Press, 1994.

\bibitem{park1995}
F.~C.~Park, J.~E.~Bobrow, and S.~R.~Ploen, ``A Lie group formulation of robot dynamics,'' \emph{The International Journal of Robotics Research}, vol.~14, no.~6, pp.~609--618, 1995.

\bibitem{bullo1995proportional}
F.~Bullo and R.~M.~Murray, ``Proportional derivative (PD) control on the Euclidean group,'' in \emph{Proc. European Control Conference}, 1995, pp.~1091--1097.

\bibitem{lee2010geometric}
T.~Lee, M.~Leok, and N.~H.~McClamroch, ``Geometric tracking control of a quadrotor UAV on SE(3),'' in \emph{Proc.\ IEEE Conference on Decision and Control}, 2010, pp.~5420--5425.

\bibitem{mellinger2011minimum}
D.~Mellinger and V.~Kumar, ``Minimum snap trajectory generation and control for quadrotors,'' in \emph{Proc. IEEE International Conference on Robotics and Automation}, 2011, pp.~2520--2525.

\bibitem{diehl2005real}
M.~Diehl, H.~G.~Bock, and J.~P.~Schl\"oder, ``A real-time iteration scheme for nonlinear optimization in optimal feedback control,'' \emph{SIAM Journal on Control and Optimization}, vol.~43, no.~5, pp.~1714--1736, 2005.

\bibitem{baumgarte1972}
J.~Baumgarte, ``Stabilization of constraints and integrals of motion in dynamical systems,'' \emph{Computer Methods in Applied Mechanics and Engineering}, vol.~1, no.~1, pp.~1--16, 1972.

\bibitem{akiba2019optuna}
T.~Akiba, S.~Sano, T.~Yanase, T.~Ohta, and M.~Koyama,
``Optuna: A next-generation hyperparameter optimization framework,''
in \emph{Proc. 25th ACM SIGKDD International Conference on
Knowledge Discovery \& Data Mining}, 2019, pp.~2623--2631.

\end{thebibliography}




\vfill

\end{document}
